If a string in $\P(A)$ has no 1-box on it, it is sent to the same
string labelled with 1. One may check that the kernel of $\Phi$ is
precisely the ideal generated by our relations $R$, so
$P^A=\P(A)/\cal J(R)$.

Observe how $P^A_n$ is a sum, over Temperley-Lieb basis tangles, of
tensor powers of $A$. When $n\!=\!2$ this gives an associative
algebra structure on $A\otimes A\oplus A\otimes A$. Explicitly,
write $(a\otimes b)\oplus 0$ as $a\otimes b$ and $0\oplus (x\otimes y)$
as $x\otimes y$. Multiplication is then determined by the rules:
\begin{eqnarray*}
(a_1\otimes b_1)(a_2\otimes b_2) & = a_1a_2\otimes b_2b_1 \\
(x_1\otimes y_1)(x_2\otimes y_2) & = {\op{tr}}(y_1x_2)x_1\otimes
y_2 \\
(a\otimes b)(x\otimes y) & = 0\oplus axb\otimes y \\
(x\otimes y)(a\otimes b) & = 0\oplus x\otimes bya
\end{eqnarray*}

The planar algebra $P^A$ may be degenerate,
even when tr on $A$ is non-degenerate and $\d$ is such that
$TL(\d)$ is non-degenerate. We will give more details on the
structure of $P^A$ in $\S$3.1.


\bigskip
{\bf Example 2.3: The Fuss-Catalan algebras} (see [BJ2]).
If $a_1,a_2,\dots ,a_k\in K-\{0\}$, $FC(n,a_1,\dots ,a_k)$ is the
algebra having as basis the Temperley-Lieb diagrams in $TL(nk)$ for
which, for each $p=1,2,\dots ,k$, the set of all boundary points
(counting from the left) indexed by
$\{jk+(-1)^jp+(\sin^2\frac{j\pi}{2})(k+j)\mid j=01,2,\dots (n-1)\}$
are connected among themselves. Assign a colour to each $p=1,2,\dots
,k$ so we think of the Temperley-Lieb strings as being coloured.
Then multiplication preserves colours so that closed loops will have
colours. Removing a closed loop coloured $m$ contributes a
multiplicative factor $a_m$. To see that $FC(n,a_1,\dots ,a_k)$ is
a planar algebra, begin with the case $k=2$. We claim $FC(n,a,b)$
is planar on one 2-box. We draw the box symbolically as
\[
        \begin{picture}(0,0)%
\epsfig{file=ps/pic16.ps}%
\end{picture}%
\setlength{\unitlength}{0.00041700in}%
%
\begingroup\makeatletter\ifx\SetFigFont\undefined
% extract first six characters in \fmtname
\def\x#1#2#3#4#5#6#7\relax{\def\x{#1#2#3#4#5#6}}%
\expandafter\x\fmtname xxxxxx\relax \def\y{splain}%
\ifx\x\y   % LaTeX or SliTeX?
\gdef\SetFigFont#1#2#3{%
  \ifnum #1<17\tiny\else \ifnum #1<20\small\else
  \ifnum #1<24\normalsize\else \ifnum #1<29\large\else
  \ifnum #1<34\Large\else \ifnum #1<41\LARGE\else
     \huge\fi\fi\fi\fi\fi\fi
  \csname #3\endcsname}%
\else
\gdef\SetFigFont#1#2#3{\begingroup
  \count@#1\relax \ifnum 25<\count@\count@25\fi
  \def\x{\endgroup\@setsize\SetFigFont{#2pt}}%
  \expandafter\x
    \csname \romannumeral\the\count@ pt\expandafter\endcsname
    \csname @\romannumeral\the\count@ pt\endcsname
  \csname #3\endcsname}%
\fi
\fi\endgroup
\begin{picture}(1544,1980)(2379,-2551)
\end{picture}

\]
This shows in fact how to define the corresponding 
$\Phi: {\cal P}_n\to FC(n,a,b)$: double all the strings and replace all the
2-boxes according to the diagram. Thus for instance the tangle 
$\cal N$ below (with boxes shrunk to points, there being only one
2-box),
\[
        \begin{picture}(0,0)%
\epsfig{file=ps/pic17.ps}%
\end{picture}%
\setlength{\unitlength}{0.00041700in}%
%
\begingroup\makeatletter\ifx\SetFigFont\undefined
% extract first six characters in \fmtname
\def\x#1#2#3#4#5#6#7\relax{\def\x{#1#2#3#4#5#6}}%
\expandafter\x\fmtname xxxxxx\relax \def\y{splain}%
\ifx\x\y   % LaTeX or SliTeX?
\gdef\SetFigFont#1#2#3{%
  \ifnum #1<17\tiny\else \ifnum #1<20\small\else
  \ifnum #1<24\normalsize\else \ifnum #1<29\large\else
  \ifnum #1<34\Large\else \ifnum #1<41\LARGE\else
     \huge\fi\fi\fi\fi\fi\fi
  \csname #3\endcsname}%
\else
\gdef\SetFigFont#1#2#3{\begingroup
  \count@#1\relax \ifnum 25<\count@\count@25\fi
  \def\x{\endgroup\@setsize\SetFigFont{#2pt}}%
  \expandafter\x
    \csname \romannumeral\the\count@ pt\expandafter\endcsname
    \csname @\romannumeral\the\count@ pt\endcsname
  \csname #3\endcsname}%
\fi
\fi\endgroup
\begin{picture}(4019,2744)(1779,-2783)
\end{picture}

\]
is sent to $\Phi(\cal N)$ below
\[
        \begin{picture}(0,0)%
\epsfig{file=ps/pic18.ps}%
\end{picture}%
\setlength{\unitlength}{0.00041700in}%
%
\begingroup\makeatletter\ifx\SetFigFont\undefined
% extract first six characters in \fmtname
\def\x#1#2#3#4#5#6#7\relax{\def\x{#1#2#3#4#5#6}}%
\expandafter\x\fmtname xxxxxx\relax \def\y{splain}%
\ifx\x\y   % LaTeX or SliTeX?
\gdef\SetFigFont#1#2#3{%
  \ifnum #1<17\tiny\else \ifnum #1<20\small\else
  \ifnum #1<24\normalsize\else \ifnum #1<29\large\else
  \ifnum #1<34\Large\else \ifnum #1<41\LARGE\else
     \huge\fi\fi\fi\fi\fi\fi
  \csname #3\endcsname}%
\else
\gdef\SetFigFont#1#2#3{\begingroup
  \count@#1\relax \ifnum 25<\count@\count@25\fi
  \def\x{\endgroup\@setsize\SetFigFont{#2pt}}%
  \expandafter\x
    \csname \romannumeral\the\count@ pt\expandafter\endcsname
    \csname @\romannumeral\the\count@ pt\endcsname
  \csname #3\endcsname}%
\fi
\fi\endgroup
\begin{picture}(8826,4142)(1779,-3801)
\put(2402,-3737){\makebox(0,0)[lb]{\smash{\SetFigFont{12}{14.4}{rm}$b$}}}
\put(3677,-3737){\makebox(0,0)[lb]{\smash{\SetFigFont{12}{14.4}{rm}$b$}}}
\put(5177,-3737){\makebox(0,0)[lb]{\smash{\SetFigFont{12}{14.4}{rm}$b$}}}
\put(2027,-3737){\makebox(0,0)[lb]{\smash{\SetFigFont{12}{14.4}{rm}$a$}}}
\put(3977,-3737){\makebox(0,0)[lb]{\smash{\SetFigFont{12}{14.4}{rm}$a$}}}
\put(4877,-3737){\makebox(0,0)[lb]{\smash{\SetFigFont{12}{14.4}{rm}$a$}}}
\put(10053,-3738){\makebox(0,0)[lb]{\smash{\SetFigFont{12}{14.4}{rm}$b$}}}
\put(9003,-3738){\makebox(0,0)[lb]{\smash{\SetFigFont{12}{14.4}{rm}$b$}}}
\put(8553,-3738){\makebox(0,0)[lb]{\smash{\SetFigFont{12}{14.4}{rm}$b$}}}
\put(8253,-3738){\makebox(0,0)[lb]{\smash{\SetFigFont{12}{14.4}{rm}$a$}}}
\put(9378,-3738){\makebox(0,0)[lb]{\smash{\SetFigFont{12}{14.4}{rm}$a$}}}
\put(9753,-3738){\makebox(0,0)[lb]{\smash{\SetFigFont{12}{14.4}{rm}$a$}}}
\put(8252, 88){\makebox(0,0)[lb]{\smash{\SetFigFont{12}{14.4}{rm}$a$}}}
\put(8552, 88){\makebox(0,0)[lb]{\smash{\SetFigFont{12}{14.4}{rm}$b$}}}
\put(9002, 88){\makebox(0,0)[lb]{\smash{\SetFigFont{12}{14.4}{rm}$b$}}}
\put(10052, 88){\makebox(0,0)[lb]{\smash{\SetFigFont{12}{14.4}{rm}$b$}}}
\put(9752, 88){\makebox(0,0)[lb]{\smash{\SetFigFont{12}{14.4}{rm}$a$}}}
\put(9377, 88){\makebox(0,0)[lb]{\smash{\SetFigFont{12}{14.4}{rm}$b$}}}
\put(5926,-1936){\makebox(0,0)[lb]{\smash{\SetFigFont{12}{14.4}{rm}$=a^2b^2$}}}
\put(5176, 89){\makebox(0,0)[lb]{\smash{\SetFigFont{12}{14.4}{rm}$b$}}}
\put(4051, 89){\makebox(0,0)[lb]{\smash{\SetFigFont{12}{14.4}{rm}$a$}}}
\put(4801, 89){\makebox(0,0)[lb]{\smash{\SetFigFont{12}{14.4}{rm}$a$}}}
\put(3676, 89){\makebox(0,0)[lb]{\smash{\SetFigFont{12}{14.4}{rm}$b$}}}
\put(2326, 89){\makebox(0,0)[lb]{\smash{\SetFigFont{12}{14.4}{rm}$a$}}}
\put(2701, 89){\makebox(0,0)[lb]{\smash{\SetFigFont{12}{14.4}{rm}$b$}}}
\end{picture}

\]
\noindent It is clear that $\Phi$ defines an algebra homomorphism
and surjectivity follows from [BJ2]. That ker
$\Phi$ is annular invariant is straightforward.
The general case of $FC(n,a_1,\dots ,a_k)$ is similar.
One considers the $k-1$ 2-boxes drawn symbolically as
\[
        \begin{picture}(0,0)%
\epsfig{file=ps/pic19.ps}%
\end{picture}%
\setlength{\unitlength}{0.00041700in}%
%
\begingroup\makeatletter\ifx\SetFigFont\undefined
% extract first six characters in \fmtname
\def\x#1#2#3#4#5#6#7\relax{\def\x{#1#2#3#4#5#6}}%
\expandafter\x\fmtname xxxxxx\relax \def\y{splain}%
\ifx\x\y   % LaTeX or SliTeX?
\gdef\SetFigFont#1#2#3{%
  \ifnum #1<17\tiny\else \ifnum #1<20\small\else
  \ifnum #1<24\normalsize\else \ifnum #1<29\large\else
  \ifnum #1<34\Large\else \ifnum #1<41\LARGE\else
     \huge\fi\fi\fi\fi\fi\fi
  \csname #3\endcsname}%
\else
\gdef\SetFigFont#1#2#3{\begingroup
  \count@#1\relax \ifnum 25<\count@\count@25\fi
  \def\x{\endgroup\@setsize\SetFigFont{#2pt}}%
  \expandafter\x
    \csname \romannumeral\the\count@ pt\expandafter\endcsname
    \csname @\romannumeral\the\count@ pt\endcsname
  \csname #3\endcsname}%
\fi
\fi\endgroup
\begin{picture}(9644,2430)(1779,-2776)
\end{picture}

\]

\noindent One proceeds as above, replacing the single strings in an
$\cal N$ by $k$ coloured strings. Surjectivity follows from
[La].

Note that these planar algebras give
invariants of systems of immersed curves with generic
singularities, and/or planar graphs. The most general such
invariant may be obtained by introducing a single 2-box which is a
linear combination of the $k-1$\quad 2-boxes described above. This will
generalize the dichromatic polynomial.

\bigskip{\bf Example 2.4: The BMW algebra.}
Let $L=L_2=\{R,Q\}$ and define the planar algebra $\cal B\cal M\cal
W$ on $L$ by the relations
%\[
%        \input{xfig/pic20}
%\]
%%%%%%%%%%%%%%%%%%%%%%%%%%%%%%%%%%%%%%%%%%%

        \begin{picture}(0,0)%
\epsfig{file=ps/pic201.ps}%
\end{picture}%
\setlength{\unitlength}{0.00041700in}%
%
\begingroup\makeatletter\ifx\SetFigFont\undefined
% extract first six characters in \fmtname
\def\x#1#2#3#4#5#6#7\relax{\def\x{#1#2#3#4#5#6}}%
\expandafter\x\fmtname xxxxxx\relax \def\y{splain}%
\ifx\x\y   % LaTeX or SliTeX?
\gdef\SetFigFont#1#2#3{%
  \ifnum #1<17\tiny\else \ifnum #1<20\small\else
  \ifnum #1<24\normalsize\else \ifnum #1<29\large\else
  \ifnum #1<34\Large\else \ifnum #1<41\LARGE\else
     \huge\fi\fi\fi\fi\fi\fi
  \csname #3\endcsname}%
\else
\gdef\SetFigFont#1#2#3{\begingroup
  \count@#1\relax \ifnum 25<\count@\count@25\fi
  \def\x{\endgroup\@setsize\SetFigFont{#2pt}}%
  \expandafter\x
    \csname \romannumeral\the\count@ pt\expandafter\endcsname
    \csname @\romannumeral\the\count@ pt\endcsname
  \csname #3\endcsname}%
\fi
\fi\endgroup
\begin{picture}(9113,2326)(2296,-1867)
\put(4097,-765){\makebox(0,0)[lb]{\smash{\SetFigFont{12}{14.4}{rm}$=$}}}
\put(6497,-768){\makebox(0,0)[lb]{\smash{\SetFigFont{12}{14.4}{rm}$=$}}}
\put(2296,-766){\makebox(0,0)[lb]{\smash{\SetFigFont{12}{14.4}{rm}$(i)$}}}
\put(5326,-841){\makebox(0,0)[lb]{\smash{\SetFigFont{12}{14.4}{rm},}}}
\put(7711,-811){\makebox(0,0)[lb]{\smash{\SetFigFont{12}{14.4}{rm},}}}
\put(9040,-810){\makebox(0,0)[lb]{\smash{\SetFigFont{12}{14.4}{rm}$=$}}}
\put(10270,-795){\makebox(0,0)[lb]{\smash{\SetFigFont{12}{14.4}{rm}$=$}}}
\put(5057,-435){\makebox(0,0)[lb]{\smash{\special{ps:gsave currentpoint currentpoint translate
-180.0 rotate neg exch neg exch translate}\SetFigFont{12}{14.4}{rm}$R$\special{ps:currentpoint grestore moveto}}}}
\put(3467,-675){\makebox(0,0)[lb]{\smash{\SetFigFont{12}{14.4}{rm}$R$}}}
\put(5866,-661){\makebox(0,0)[lb]{\smash{\SetFigFont{12}{14.4}{rm}$Q$}}}
\put(7471,-406){\makebox(0,0)[lb]{\smash{\special{ps:gsave currentpoint currentpoint translate
-180.0 rotate neg exch neg exch translate}\SetFigFont{12}{14.4}{rm}$Q$\special{ps:currentpoint grestore moveto}}}}
\put(8432,-180){\makebox(0,0)[lb]{\smash{\SetFigFont{12}{14.4}{rm}$R$}}}
\put(8446,-1231){\makebox(0,0)[lb]{\smash{\SetFigFont{12}{14.4}{rm}$Q$}}}
\put(10996,-211){\makebox(0,0)[lb]{\smash{\SetFigFont{12}{14.4}{rm}$Q$}}}
\put(10982,-1290){\makebox(0,0)[lb]{\smash{\SetFigFont{12}{14.4}{rm}$R$}}}
\end{picture}



        \begin{picture}(0,0)%
\epsfig{file=ps/pic202.ps}%
\end{picture}%
\setlength{\unitlength}{0.00041700in}%
%
\begingroup\makeatletter\ifx\SetFigFont\undefined
% extract first six characters in \fmtname
\def\x#1#2#3#4#5#6#7\relax{\def\x{#1#2#3#4#5#6}}%
\expandafter\x\fmtname xxxxxx\relax \def\y{splain}%
\ifx\x\y   % LaTeX or SliTeX?
\gdef\SetFigFont#1#2#3{%
  \ifnum #1<17\tiny\else \ifnum #1<20\small\else
  \ifnum #1<24\normalsize\else \ifnum #1<29\large\else
  \ifnum #1<34\Large\else \ifnum #1<41\LARGE\else
     \huge\fi\fi\fi\fi\fi\fi
  \csname #3\endcsname}%
\else
\gdef\SetFigFont#1#2#3{\begingroup
  \count@#1\relax \ifnum 25<\count@\count@25\fi
  \def\x{\endgroup\@setsize\SetFigFont{#2pt}}%
  \expandafter\x
    \csname \romannumeral\the\count@ pt\expandafter\endcsname
    \csname @\romannumeral\the\count@ pt\endcsname
  \csname #3\endcsname}%
\fi
\fi\endgroup
\begin{picture}(8785,1905)(2281,-3969)
\put(2281,-3106){\makebox(0,0)[lb]{\smash{\SetFigFont{12}{14.4}{rm}$(ii)$}}}
\put(6250,-3165){\makebox(0,0)[lb]{\smash{\SetFigFont{12}{14.4}{rm}$=$}}}
\put(8290,-3165){\makebox(0,0)[lb]{\smash{\SetFigFont{12}{14.4}{rm}$=$}}}
\put(4067,-3090){\makebox(0,0)[lb]{\smash{\SetFigFont{12}{14.4}{rm}$R$}}}
\put(5626,-3091){\makebox(0,0)[lb]{\smash{\SetFigFont{12}{14.4}{rm}$Q$}}}
\put(9106,-3061){\makebox(0,0)[lb]{\smash{\SetFigFont{12}{14.4}{rm}$Q$}}}
\put(10699,-3076){\makebox(0,0)[lb]{\smash{\SetFigFont{12}{14.4}{rm}$R$}}}
\end{picture}



        \begin{picture}(0,0)%
\epsfig{file=ps/pic203.ps}%
\end{picture}%
\setlength{\unitlength}{0.00041700in}%
%
\begingroup\makeatletter\ifx\SetFigFont\undefined
% extract first six characters in \fmtname
\def\x#1#2#3#4#5#6#7\relax{\def\x{#1#2#3#4#5#6}}%
\expandafter\x\fmtname xxxxxx\relax \def\y{splain}%
\ifx\x\y   % LaTeX or SliTeX?
\gdef\SetFigFont#1#2#3{%
  \ifnum #1<17\tiny\else \ifnum #1<20\small\else
  \ifnum #1<24\normalsize\else \ifnum #1<29\large\else
  \ifnum #1<34\Large\else \ifnum #1<41\LARGE\else
     \huge\fi\fi\fi\fi\fi\fi
  \csname #3\endcsname}%
\else
\gdef\SetFigFont#1#2#3{\begingroup
  \count@#1\relax \ifnum 25<\count@\count@25\fi
  \def\x{\endgroup\@setsize\SetFigFont{#2pt}}%
  \expandafter\x
    \csname \romannumeral\the\count@ pt\expandafter\endcsname
    \csname @\romannumeral\the\count@ pt\endcsname
  \csname #3\endcsname}%
\fi
\fi\endgroup
\begin{picture}(12138,1904)(2266,-7063)
\put(2266,-6121){\makebox(0,0)[lb]{\smash{\SetFigFont{12}{14.4}{rm}$(iii)$}}}
\put(10816,-6129){\makebox(0,0)[lb]{\smash{\SetFigFont{12}{14.4}{rm}$=$}}}
\put(8494,-6253){\makebox(0,0)[lb]{\smash{\SetFigFont{12}{14.4}{rm},}}}
\put(7471,-6121){\makebox(0,0)[lb]{\smash{\SetFigFont{12}{14.4}{rm}$Q$}}}
\put(9556,-6121){\makebox(0,0)[lb]{\smash{\SetFigFont{12}{14.4}{rm}$Q$}}}
\put(12904,-6106){\makebox(0,0)[lb]{\smash{\SetFigFont{12}{14.4}{rm}$R$}}}
\put(12211,-6159){\makebox(0,0)[lb]{\smash{\SetFigFont{12}{14.4}{rm}$=$}}}
\put(11167,-6148){\makebox(0,0)[lb]{\smash{\SetFigFont{12}{14.4}{rm}$a^{-1}$}}}
\put(5425,-6150){\makebox(0,0)[lb]{\smash{\SetFigFont{12}{14.4}{rm}$=$}}}
\put(6076,-6184){\makebox(0,0)[lb]{\smash{\SetFigFont{12}{14.4}{rm}$a$}}}
\put(6775,-6165){\makebox(0,0)[lb]{\smash{\SetFigFont{12}{14.4}{rm}$=$}}}
\put(4129,-6151){\makebox(0,0)[lb]{\smash{\SetFigFont{12}{14.4}{rm}$R$}}}
\put(13984,-6223){\makebox(0,0)[lb]{\smash{\SetFigFont{12}{14.4}{rm},}}}
\put(14404,-6193){\makebox(0,0)[lb]{\smash{\SetFigFont{12}{14.4}{rm}$(a\in C-\{0\})$}}}
\end{picture}



        \begin{picture}(0,0)%
\epsfig{file=ps/pic204.ps}%
\end{picture}%
\setlength{\unitlength}{0.00041700in}%
%
\begingroup\makeatletter\ifx\SetFigFont\undefined
% extract first six characters in \fmtname
\def\x#1#2#3#4#5#6#7\relax{\def\x{#1#2#3#4#5#6}}%
\expandafter\x\fmtname xxxxxx\relax \def\y{splain}%
\ifx\x\y   % LaTeX or SliTeX?
\gdef\SetFigFont#1#2#3{%
  \ifnum #1<17\tiny\else \ifnum #1<20\small\else
  \ifnum #1<24\normalsize\else \ifnum #1<29\large\else
  \ifnum #1<34\Large\else \ifnum #1<41\LARGE\else
     \huge\fi\fi\fi\fi\fi\fi
  \csname #3\endcsname}%
\else
\gdef\SetFigFont#1#2#3{\begingroup
  \count@#1\relax \ifnum 25<\count@\count@25\fi
  \def\x{\endgroup\@setsize\SetFigFont{#2pt}}%
  \expandafter\x
    \csname \romannumeral\the\count@ pt\expandafter\endcsname
    \csname @\romannumeral\the\count@ pt\endcsname
  \csname #3\endcsname}%
\fi
\fi\endgroup
\begin{picture}(7756,3100)(2267,-2575)
\put(2267,-1142){\makebox(0,0)[lb]{\smash{\SetFigFont{12}{14.4}{rm}$(iv)$}}}
\put(6827,-1187){\makebox(0,0)[lb]{\smash{\SetFigFont{12}{14.4}{rm}$=$}}}
\put(4187,-467){\makebox(0,0)[lb]{\smash{\SetFigFont{12}{14.4}{rm}$Q$}}}
\put(6062,-1202){\makebox(0,0)[lb]{\smash{\special{ps:gsave currentpoint currentpoint translate
-90.0 rotate neg exch neg exch translate}\SetFigFont{12}{14.4}{rm}$Q$\special{ps:currentpoint grestore moveto}}}}
\put(9152,-122){\makebox(0,0)[lb]{\smash{\special{ps:gsave currentpoint currentpoint translate
-90.0 rotate neg exch neg exch translate}\SetFigFont{12}{14.4}{rm}$Q$\special{ps:currentpoint grestore moveto}}}}
\put(9152,-2342){\makebox(0,0)[lb]{\smash{\special{ps:gsave currentpoint currentpoint translate
-90.0 rotate neg exch neg exch translate}\SetFigFont{12}{14.4}{rm}$R$\special{ps:currentpoint grestore moveto}}}}
\put(4202,-2012){\makebox(0,0)[lb]{\smash{\SetFigFont{12}{14.4}{rm}$R$}}}
\put(7757,-1187){\makebox(0,0)[lb]{\smash{\SetFigFont{12}{14.4}{rm}$R$}}}
\end{picture}



        \begin{picture}(0,0)%
\epsfig{file=ps/pic205.ps}%
\end{picture}%
\setlength{\unitlength}{0.00041700in}%
%
\begingroup\makeatletter\ifx\SetFigFont\undefined
% extract first six characters in \fmtname
\def\x#1#2#3#4#5#6#7\relax{\def\x{#1#2#3#4#5#6}}%
\expandafter\x\fmtname xxxxxx\relax \def\y{splain}%
\ifx\x\y   % LaTeX or SliTeX?
\gdef\SetFigFont#1#2#3{%
  \ifnum #1<17\tiny\else \ifnum #1<20\small\else
  \ifnum #1<24\normalsize\else \ifnum #1<29\large\else
  \ifnum #1<34\Large\else \ifnum #1<41\LARGE\else
     \huge\fi\fi\fi\fi\fi\fi
  \csname #3\endcsname}%
\else
\gdef\SetFigFont#1#2#3{\begingroup
  \count@#1\relax \ifnum 25<\count@\count@25\fi
  \def\x{\endgroup\@setsize\SetFigFont{#2pt}}%
  \expandafter\x
    \csname \romannumeral\the\count@ pt\expandafter\endcsname
    \csname @\romannumeral\the\count@ pt\endcsname
  \csname #3\endcsname}%
\fi
\fi\endgroup
\begin{picture}(8615,1725)(2297,-1374)
\put(4847,-528){\makebox(0,0)[lb]{\smash{\SetFigFont{12}{14.4}{rm}$+$}}}
\put(8826,-714){\makebox(0,0)[lb]{\smash{\SetFigFont{12}{14.4}{rm}$+$}}}
\put(2297,-543){\makebox(0,0)[lb]{\smash{\SetFigFont{12}{14.4}{rm}$(v)$}}}
\put(4157,-452){\makebox(0,0)[lb]{\smash{\SetFigFont{12}{14.4}{rm}$R$}}}
\put(5702,-452){\makebox(0,0)[lb]{\smash{\SetFigFont{12}{14.4}{rm}$Q$}}}
\put(6407,-558){\makebox(0,0)[lb]{\smash{\SetFigFont{12}{14.4}{rm}$=$}}}
\put(7235,-584){\makebox(0,0)[lb]{\smash{\SetFigFont{12}{14.4}{rm}$x$}}}
\end{picture}


%%%%%%%%%%%%%%%%%%%%%%%%%%%%%%%%%%%%%%%%%%
Note that we could use relation (v) to express $\cal B\cal M\cal
W$ using only the one label $R$, but the relations would then be
more complicated. At this stage $\cal B\cal M\cal
W$ could be zero or infinite dimensional, but we may define a homomorphism from $\cal B\cal M\cal
W$ to the algebra $BMW$ of [BiW],[Mu] by sending \
$\blankbox{R}$ to $\crss$ and  $\blankbox{Q}$ to $\crssr$ .
This homomorphism is obviously surjective and one may use the
dimension count of [BiW] to show also that
dim ${\cal BMW}(n)\leq 1.3.5.\dots(2n-1)$ so that
${\cal BMW}\cong BMW$ as algebras. Thus $BMW$ is planar.
It is also connected and the invariant of planar networks is the
Kauffman regular-isotopy two-variable polynomial of [Ka2].

\bigskip{\bf Remark.} Had we presented $BMW$ on the single 2-box
$R$, the Reidemeister type III move (number (iv) above) would have
been
\[
        \begin{picture}(0,0)%
\epsfig{file=ps/pic21.ps}%
\end{picture}%
\setlength{\unitlength}{0.00041700in}%
%
\begingroup\makeatletter\ifx\SetFigFont\undefined
% extract first six characters in \fmtname
\def\x#1#2#3#4#5#6#7\relax{\def\x{#1#2#3#4#5#6}}%
\expandafter\x\fmtname xxxxxx\relax \def\y{splain}%
\ifx\x\y   % LaTeX or SliTeX?
\gdef\SetFigFont#1#2#3{%
  \ifnum #1<17\tiny\else \ifnum #1<20\small\else
  \ifnum #1<24\normalsize\else \ifnum #1<29\large\else
  \ifnum #1<34\Large\else \ifnum #1<41\LARGE\else
     \huge\fi\fi\fi\fi\fi\fi
  \csname #3\endcsname}%
\else
\gdef\SetFigFont#1#2#3{\begingroup
  \count@#1\relax \ifnum 25<\count@\count@25\fi
  \def\x{\endgroup\@setsize\SetFigFont{#2pt}}%
  \expandafter\x
    \csname \romannumeral\the\count@ pt\expandafter\endcsname
    \csname @\romannumeral\the\count@ pt\endcsname
  \csname #3\endcsname}%
\fi
\fi\endgroup
\begin{picture}(10765,4050)(2341,-3391)
\put(2746,-2236){\makebox(0,0)[lb]{\smash{\SetFigFont{12}{14.4}{rm}$R$}}}
\put(5881,-1396){\makebox(0,0)[lb]{\smash{\SetFigFont{12}{14.4}{rm}$=$}}}
\put(8311,-2564){\makebox(0,0)[lb]{\smash{\special{ps:gsave currentpoint currentpoint translate
-90.0 rotate neg exch neg exch translate}\SetFigFont{12}{14.4}{rm}$R$\special{ps:currentpoint grestore moveto}}}}
\put(6916,-1409){\makebox(0,0)[lb]{\smash{\SetFigFont{12}{14.4}{rm}$R$}}}
\put(8326,-344){\makebox(0,0)[lb]{\smash{\special{ps:gsave currentpoint currentpoint translate
-90.0 rotate neg exch neg exch translate}\SetFigFont{12}{14.4}{rm}$R$\special{ps:currentpoint grestore moveto}}}}
\put(4606,-1469){\makebox(0,0)[lb]{\smash{\special{ps:gsave currentpoint currentpoint translate
-90.0 rotate neg exch neg exch translate}\SetFigFont{12}{14.4}{rm}$R$\special{ps:currentpoint grestore moveto}}}}
\put(2716,-676){\makebox(0,0)[lb]{\smash{\SetFigFont{12}{14.4}{rm}$R$}}}
\put(9351,-1396){\makebox(0,0)[lb]{\smash{\SetFigFont{12}{14.4}{rm}$+$ }}}
\put(10101,-1171){\makebox(0,0)[lb]{\smash{\SetFigFont{12}{14.4}{rm}pictures with }}}
\put(10101,-1675){\makebox(0,0)[lb]{\smash{\SetFigFont{12}{14.4}{rm}at most two}}}
\put(12506,-1606){\makebox(0,0)[lb]{\smash{\SetFigFont{12}{14.4}{rm}$R$}}}
\put(13106,-1696){\makebox(0,0)[lb]{\smash{\SetFigFont{12}{14.4}{rm}'s}}}
\end{picture}

\]
\noindent
This leads us to consider the general planar algebra $B_n$ with the
following three conditions:
\begin{verse}
(1) $B_n$ is planar on one 2-box.

(2) dim $B_2=3$.

(3) dim $B_3\leq 15$.
\end{verse}
If one lists 16 tangles in $B_3$ then generically any one of them
will have to be a linear combination of the other $15$. Looking at
the 15th and 16th tangles in a listing according to the number of
2-boxes occuring in the tangle, we will generically obtain a type III Reidemeister 
move, or Yang-Baxter equation, modulo terms with less 2-boxes,as above.
It is not hard to show that these conditions force
dim $B_n\leq 1.3.5\cdot\dots \cdot (2n-1)$ since there are
necessarily Reidemeister-like moves of types I and II.
Note that $FC(n,a,b)$ satisfies these conditions as well as $BMW$!
For $C^*$-planar algebras, we have shown with Bisch
([BJ1]) that the only $B$'s with (1) and (2) as above,
and dim $B_3\leq 12$ are the Fuss-Catalan algebras  (with one
exception, when dim $B_3=9$).

\bigskip{\bf Example 2.5: a Hecke-algebra related example.}
The {\sc homfly} polynomial of [F+] is highly sensitive to the
orientation of a link and we may not proceed to use it to define a
planar algebra as in Example 2.4. In particular, a crossing in the
{\sc homfly} theory is necessarily oriented as $\crsso$.
Thus it does {\it not} yield a 2-box in our planar algebra context.
Nevertheless it is possible to use the {\sc homfly} skein theory to
define a planar algebra. We let $P^H_k$ be the usual {\sc homfly}
skein algebra of linear combinations of (3-dimensional) isotopy
classes of oriented tangles in the product of the $k$-box with an
interval, with orientations alternating out-in, modulo the
{\sc homfly} skein relation\newline
 $t\crsso -t^{-1}\crssor =x\smthv$ 
where $t\neq 0$ and $x$ are scalars. Projected onto the $k$-box,
such a tangle could look as in Figure 2.5.1,

\[
	\begin{picture}(0,0)%
\epsfig{file=ps/pic22.ps}%
\end{picture}%
\setlength{\unitlength}{0.00041700in}%
%
\begingroup\makeatletter\ifx\SetFigFont\undefined
% extract first six characters in \fmtname
\def\x#1#2#3#4#5#6#7\relax{\def\x{#1#2#3#4#5#6}}%
\expandafter\x\fmtname xxxxxx\relax \def\y{splain}%
\ifx\x\y   % LaTeX or SliTeX?
\gdef\SetFigFont#1#2#3{%
  \ifnum #1<17\tiny\else \ifnum #1<20\small\else
  \ifnum #1<24\normalsize\else \ifnum #1<29\large\else
  \ifnum #1<34\Large\else \ifnum #1<41\LARGE\else
     \huge\fi\fi\fi\fi\fi\fi
  \csname #3\endcsname}%
\else
\gdef\SetFigFont#1#2#3{\begingroup
  \count@#1\relax \ifnum 25<\count@\count@25\fi
  \def\x{\endgroup\@setsize\SetFigFont{#2pt}}%
  \expandafter\x
    \csname \romannumeral\the\count@ pt\expandafter\endcsname
    \csname @\romannumeral\the\count@ pt\endcsname
  \csname #3\endcsname}%
\fi
\fi\endgroup
\begin{picture}(4281,3435)(553,-3751)
\end{picture}

\]
\begin{center}
Figure 2.5.1
\end{center}

If we take the labeling set $L_k=P^H_k$, then $P^H$ is a general planar
algebra since planar isotopy of projections implies 3-dimensional
isotopy (invariance under the annular category is easy).
Standard {\sc homfly} arguments show dim $(P^H_k)\leq k!$
and a specialization could be used to obtain equality. Thus the
algebra is planar and the invariant is clearly the
{\sc homfly} polynomial of the
oriented link diagram given by a labeled network in $P^H_0$.
If we used this invariant to define the algebra as in $\S$1, we
would only obtain the same algebra for generic values of $t$ and
$x$. Note that $P^H_k$ is {\it not} isomorphic to the Hecke algebra
for $k\geq 4$, e.g. $P^H_4$ has an irreducible 4-dimensional
representation. In fact $P^H_k$ is, for generic $(t,x)$ and large
$n$, isomorphic to End${}_{SU(n)}
\underbrace{(V\otimes\bar V\otimes V\otimes \bar V\dots)}_{k\text{ \ vector\  spaces}}$,
where $V=\Bbb C^n$, the obvious $SU(n)$-module. This isomorphism is
only an algebra isomorphism, not a planar algebra isomorphism.

It is clear that the labeling set for $P^H$ could be reduced to a
set of $k!$ isotopy classes of tangles for $P^H_k$. But in fact a
single 3-label suffices as we now show

\bs
\ni{\bf Theorem 2.5.2} Any tangle {\it in the knot-theoretic sense}
with alternating in and out boundary orientations is
isotopic to a tangle with a diagram where all crossings occur in
disjoint discs which contain the pattern 
\[
	\begin{picture}(0,0)%
\epsfig{file=ps/pic23.ps}%
\end{picture}%
\setlength{\unitlength}{0.00041700in}%
%
\begingroup\makeatletter\ifx\SetFigFont\undefined
% extract first six characters in \fmtname
\def\x#1#2#3#4#5#6#7\relax{\def\x{#1#2#3#4#5#6}}%
\expandafter\x\fmtname xxxxxx\relax \def\y{splain}%
\ifx\x\y   % LaTeX or SliTeX?
\gdef\SetFigFont#1#2#3{%
  \ifnum #1<17\tiny\else \ifnum #1<20\small\else
  \ifnum #1<24\normalsize\else \ifnum #1<29\large\else
  \ifnum #1<34\Large\else \ifnum #1<41\LARGE\else
     \huge\fi\fi\fi\fi\fi\fi
  \csname #3\endcsname}%
\else
\gdef\SetFigFont#1#2#3{\begingroup
  \count@#1\relax \ifnum 25<\count@\count@25\fi
  \def\x{\endgroup\@setsize\SetFigFont{#2pt}}%
  \expandafter\x
    \csname \romannumeral\the\count@ pt\expandafter\endcsname
    \csname @\romannumeral\the\count@ pt\endcsname
  \csname #3\endcsname}%
\fi
\fi\endgroup
\begin{picture}(1230,1230)(886,-976)
\end{picture}

\]
with some non-alternating choice of crossings.

\bs
{\sc Proof.} We begin with a tangle without boundary, i.e.~an
oriented link $L$. Choose a diagram for $L$ and add a parallel
double $L'$ of $L$ to the left of $L$ and oppositely oriented, with
crossings chosen so that

\begin{itemize}
\item[a)] $L'$ is always under $L$

\item[b)] $L'$  itself is an unlink.
\end{itemize}
An example of the resulting diagram (for the Whitehead Link) is
given in Figure 2.5.3
\[
	\begin{picture}(0,0)%
\epsfig{file=ps/pic24.ps}%
\end{picture}%
\setlength{\unitlength}{0.00041700in}%
%
\begingroup\makeatletter\ifx\SetFigFont\undefined
% extract first six characters in \fmtname
\def\x#1#2#3#4#5#6#7\relax{\def\x{#1#2#3#4#5#6}}%
\expandafter\x\fmtname xxxxxx\relax \def\y{splain}%
\ifx\x\y   % LaTeX or SliTeX?
\gdef\SetFigFont#1#2#3{%
  \ifnum #1<17\tiny\else \ifnum #1<20\small\else
  \ifnum #1<24\normalsize\else \ifnum #1<29\large\else
  \ifnum #1<34\Large\else \ifnum #1<41\LARGE\else
     \huge\fi\fi\fi\fi\fi\fi
  \csname #3\endcsname}%
\else
\gdef\SetFigFont#1#2#3{\begingroup
  \count@#1\relax \ifnum 25<\count@\count@25\fi
  \def\x{\endgroup\@setsize\SetFigFont{#2pt}}%
  \expandafter\x
    \csname \romannumeral\the\count@ pt\expandafter\endcsname
    \csname @\romannumeral\the\count@ pt\endcsname
  \csname #3\endcsname}%
\fi
\fi\endgroup
\begin{picture}(4473,6470)(1684,-6412)
\end{picture}

\]
\begin{center}
	Figure 2.5.3
\end{center}
\ni Now join $L$ to $L'$, component by component, by replacing
$\paro$  by $\smthvo$, at some point well away from any crossings. 
Since $L'$ is
an unlink below $L$, the resulting link is isotopic to $L$. 
All the crossings in $L\cup L'$ occur in disjoint discs containing the pattern
\begin{picture}(0,0)%
\epsfig{file=ps/pic25.ps}%
\end{picture}%
\setlength{\unitlength}{0.00041700in}%
%
\begingroup\makeatletter\ifx\SetFigFont\undefined
% extract first six characters in \fmtname
\def\x#1#2#3#4#5#6#7\relax{\def\x{#1#2#3#4#5#6}}%
\expandafter\x\fmtname xxxxxx\relax \def\y{splain}%
\ifx\x\y   % LaTeX or SliTeX?
\gdef\SetFigFont#1#2#3{%
  \ifnum #1<17\tiny\else \ifnum #1<20\small\else
  \ifnum #1<24\normalsize\else \ifnum #1<29\large\else
  \ifnum #1<34\Large\else \ifnum #1<41\LARGE\else
     \huge\fi\fi\fi\fi\fi\fi
  \csname #3\endcsname}%
\else
\gdef\SetFigFont#1#2#3{\begingroup
  \count@#1\relax \ifnum 25<\count@\count@25\fi
  \def\x{\endgroup\@setsize\SetFigFont{#2pt}}%
  \expandafter\x
    \csname \romannumeral\the\count@ pt\expandafter\endcsname
    \csname @\romannumeral\the\count@ pt\endcsname
  \csname #3\endcsname}%
\fi
\fi\endgroup
\begin{picture}(1298,1288)(548,-1305)
\end{picture}
 
which can be isotoped to the pattern
\begin{picture}(0,0)%
\epsfig{file=ps/pic26.ps}%
\end{picture}%
\setlength{\unitlength}{0.00041700in}%
%
\begingroup\makeatletter\ifx\SetFigFont\undefined
% extract first six characters in \fmtname
\def\x#1#2#3#4#5#6#7\relax{\def\x{#1#2#3#4#5#6}}%
\expandafter\x\fmtname xxxxxx\relax \def\y{splain}%
\ifx\x\y   % LaTeX or SliTeX?
\gdef\SetFigFont#1#2#3{%
  \ifnum #1<17\tiny\else \ifnum #1<20\small\else
  \ifnum #1<24\normalsize\else \ifnum #1<29\large\else
  \ifnum #1<34\Large\else \ifnum #1<41\LARGE\else
     \huge\fi\fi\fi\fi\fi\fi
  \csname #3\endcsname}%
\else
\gdef\SetFigFont#1#2#3{\begingroup
  \count@#1\relax \ifnum 25<\count@\count@25\fi
  \def\x{\endgroup\@setsize\SetFigFont{#2pt}}%
  \expandafter\x
    \csname \romannumeral\the\count@ pt\expandafter\endcsname
    \csname @\romannumeral\the\count@ pt\endcsname
  \csname #3\endcsname}%
\fi
\fi\endgroup
\begin{picture}(1260,1289)(571,-1313)
\end{picture}
 
containing two discs of the required form. 
Alternating patterns can be
avoided by keeping the top string on top in this isotopy.

For a tangle with boundary we make the doubling curve follow the
boundary, turning right just before it would hit it and right
again as it nears the point where the next string exits the
tangle, as in Figure 2.5.4
\[
	\begin{picture}(0,0)%
\epsfig{file=ps/pic27.ps}%
\end{picture}%
\setlength{\unitlength}{0.00041700in}%
%
\begingroup\makeatletter\ifx\SetFigFont\undefined
% extract first six characters in \fmtname
\def\x#1#2#3#4#5#6#7\relax{\def\x{#1#2#3#4#5#6}}%
\expandafter\x\fmtname xxxxxx\relax \def\y{splain}%
\ifx\x\y   % LaTeX or SliTeX?
\gdef\SetFigFont#1#2#3{%
  \ifnum #1<17\tiny\else \ifnum #1<20\small\else
  \ifnum #1<24\normalsize\else \ifnum #1<29\large\else
  \ifnum #1<34\Large\else \ifnum #1<41\LARGE\else
     \huge\fi\fi\fi\fi\fi\fi
  \csname #3\endcsname}%
\else
\gdef\SetFigFont#1#2#3{\begingroup
  \count@#1\relax \ifnum 25<\count@\count@25\fi
  \def\x{\endgroup\@setsize\SetFigFont{#2pt}}%
  \expandafter\x
    \csname \romannumeral\the\count@ pt\expandafter\endcsname
    \csname @\romannumeral\the\count@ pt\endcsname
  \csname #3\endcsname}%
\fi
\fi\endgroup
\begin{picture}(5745,2926)(586,-3998)
\put(6331,-1336){\makebox(0,0)[lb]{\smash{\SetFigFont{12}{14.4}{rm}boundary of tangle}}}
\end{picture}

\]
\begin{center}
	Figure 2.5.4
\end{center}
\ni Join the original tangle to $L'$ one string at a time and
proceed as before. \qed

\bs\ni{\bf Corollary 2.5.5} The planar algebra $P^H$ is generated by
the single 3-box \begin{picture}(0,0)%
\epsfig{file=ps/pic28.ps}%
\end{picture}%
\setlength{\unitlength}{0.00041700in}%
%
\begingroup\makeatletter\ifx\SetFigFont\undefined
% extract first six characters in \fmtname
\def\x#1#2#3#4#5#6#7\relax{\def\x{#1#2#3#4#5#6}}%
\expandafter\x\fmtname xxxxxx\relax \def\y{splain}%
\ifx\x\y   % LaTeX or SliTeX?
\gdef\SetFigFont#1#2#3{%
  \ifnum #1<17\tiny\else \ifnum #1<20\small\else
  \ifnum #1<24\normalsize\else \ifnum #1<29\large\else
  \ifnum #1<34\Large\else \ifnum #1<41\LARGE\else
     \huge\fi\fi\fi\fi\fi\fi
  \csname #3\endcsname}%
\else
\gdef\SetFigFont#1#2#3{\begingroup
  \count@#1\relax \ifnum 25<\count@\count@25\fi
  \def\x{\endgroup\@setsize\SetFigFont{#2pt}}%
  \expandafter\x
    \csname \romannumeral\the\count@ pt\expandafter\endcsname
    \csname @\romannumeral\the\count@ pt\endcsname
  \csname #3\endcsname}%
\fi
\fi\endgroup
\begin{picture}(1844,1267)(879,-1883)
\end{picture}
.

\bs
{\sc Proof.} The {\sc homfly} relations can be used to go between
the various possible choices of crossings in the 3-box of Theorem
2.5.2. \qed

\bs This corollary was first proved  by W.~B.~R.~Lickorish using an
argument adapted to the {\sc homfly} skein. His argument is much
more efficient in producing a skein element involving only the
above 3-box. The tangles may be chosen
{\it alternating} in Theorem 2.5.2.

\bs{\bf Remarks.} 1) Another way of stating Theorem 2.5.2 is to say
that any tangle can be projected with only simple tiple point
singularities.  One may ask if there are a set of ``Reidemeister
moves" for such non-generic projections.

2) A related question would be to find a presentation of $P^H$ on
the above 3-box.


\bs{\bf Discussion 2.5.6}  In the remark of Example 2.4 we
introduced relations on the planar algebra generated by a 2-box,
which force finite dimensionality of all the $P_n$'s. One should
explore the possibilities for the planar algebra generated by a
single 3-box. The dimension restrictions analogous to the 1,3 $\leq
15$ values of Example 2.4 are 1,2,6,$\leq 24$ and we conjecture,
somewhat weakly, the following

\bs
\ni
{\bf Conjecture 2.5.7} Let $(P,\Phi)$ be a planar algebra with
labelling set $L=L_3$, $\#(L_3)=1$. Suppose $\dim P_n\leq n!$ for
$n\leq 4$. Let $V$ be the subspace of ${\cal P}_4(L)$ spanned by
tangles with at most two labeled 3-boxes, and let
$R=V\cap{\op{ker}} \ \Phi$ be relations. Then
\[
\dim\left( \frac{{\cal P}_n(L)}{{\cal J}_n(R)}\right)\leq n! \quad
{\text{for all}} \ n.
\]

\bigskip
There is some evidence for the conjecture. It would imply in
particular the $n=0$ case which implies the following
result, proved by D. Thurston about hexagons:

%\block
``Consider all graphs with hexagonal faces that may be drawn on
$S^2$ with non-intersecting edges. Let $M$ be the move of Figure
2.5.8 on the set of all such graphs (where the 8 external vertices
are connected in an arbitrary way to the rest of the graph).
\[
	\begin{picture}(0,0)%
\epsfig{file=ps/pic29.ps}%
\end{picture}%
\setlength{\unitlength}{0.00041700in}%
%
\begingroup\makeatletter\ifx\SetFigFont\undefined
% extract first six characters in \fmtname
\def\x#1#2#3#4#5#6#7\relax{\def\x{#1#2#3#4#5#6}}%
\expandafter\x\fmtname xxxxxx\relax \def\y{splain}%
\ifx\x\y   % LaTeX or SliTeX?
\gdef\SetFigFont#1#2#3{%
  \ifnum #1<17\tiny\else \ifnum #1<20\small\else
  \ifnum #1<24\normalsize\else \ifnum #1<29\large\else
  \ifnum #1<34\Large\else \ifnum #1<41\LARGE\else
     \huge\fi\fi\fi\fi\fi\fi
  \csname #3\endcsname}%
\else
\gdef\SetFigFont#1#2#3{\begingroup
  \count@#1\relax \ifnum 25<\count@\count@25\fi
  \def\x{\endgroup\@setsize\SetFigFont{#2pt}}%
  \expandafter\x
    \csname \romannumeral\the\count@ pt\expandafter\endcsname
    \csname @\romannumeral\the\count@ pt\endcsname
  \csname #3\endcsname}%
\fi
\fi\endgroup
\begin{picture}(7320,2940)(850,-2725)
\end{picture}

\]
\begin{center}
	Figure 2.5.8: The move $M$
\end{center}

\medskip\ni Then one may find a finite number of applications of
the move $M$ leading to a graph with two adjacent 2-valent
vertices." 
%\endblock

\bs
\bs
{\bf Example 2.6: Tensors.}
Let $V$ be a finite dimensional vector space with dual $\hat V$.
We will define a planar algebra $P^{\otimes}\!\!=\!\cup_k
P^{\otimes}_k$ with dim $P_0^{\otimes}\!=\!0$ and
$P_k^{\otimes}\!\!={\op{End}}(V\!\otimes\!\hat V\!\otimes
\!V\!\otimes\!
\hat V\!\otimes\!\dots)$ where there are $k$ vector spaces in the
tensor product. The planar structure on $P^{\otimes}$ can be defined
invariantly using the canonical maps $V\otimes\hat V\to K$ and
$\hat V\otimes  V\to K$, which are applied to any pair of vector
spaces connected by an internal edge in a planar tangle, where $V$
and $\hat V$ are associated with the marked points of a $k$-box in
an alternating fashion with $V$ associated to $*$. One could
also think of $\Bbb C\subset$End$(V)$ as a finite factor
and use the method of Theorem 4.2.1.
 It is perhaps easiest to understand this structure using
a basis
$(v_1,v_2,\dots ,v_n)$ of $V$, with corresponding dual basis.
An element of
$P^\otimes_k$ is then the same as a tensor
$X^{j_1j_2\dots j_k}_{i_1i_2\dots i_k}$. The labelling set $L_k$ is
$P^\otimes_k$ itself and the presenting map $\Phi:{\cal P}_k(L)\to\
P^\otimes_k$ is defined by summing (``contracting") over all the
internal indices  in a labelled planar tangle. The first marked
point in a box corresponds to the ``$j_1$" above. To be more
precise, one defines a {\it state} $\sigma$ of a planar tangle $T$
to be a function from the connected components of the set $S(T)$
of strings in $T$, \ $\sigma: S(T)\to\{1,2,\dots ,n\}$, to the basis
elements of $V$. A state defines a set of indices around every box
$B$ in $T$, and since the label associated to $B$ is a tensor, with
the appropriate number of indices, to each labelled box, the state
$\sigma$ associates a number, $\sigma(B)$. A state also induces a
function $\p\sigma$ from the marked points on the boundary of $T$
to $\{1,2,\dots , n\}$. Now we associate a tensor $\Phi(T)$ with
$T$ as follows: let $f:$ \{marked points $(T)\}\to\{1,2,\dots ,n\}$
denote the indices 
$\dsize{{j_1\ldots j_k}\choose {i_1 \ldots  i_k}}$
%$\pmatrix j_1\dots j_k\\i_1\dots i_k\endpmatrix$
of a tensor in $P^\otimes_k\quad (f(p,0)=i_p,f(p,1)=j_p,$ for $ 1\leq p \leq k)$. Then
$$
\Phi(T)^{j_1\dots j_k}_{i_1\dots i_k}=\sum_{\sigma:\p\sigma =f}
\left( \prod_{B\in\{ {\text{labelled}}\atop {\text{boxes of}} \ T\}
} \sigma(B)\right) \ .
$$
An empty sum is zero and an empty product is 1. One easily checks
that $\Phi$ defines an algebra homomorphism
and that ker $\Phi$ is invariant under the annular category.

This planar algebra has an obvious
$*$-structure. The invariant $Z$ is recognizable as the partition
function for the ``vertex model" defined by the labelled network,
the labels supplying the Boltzmann weights (see [Ba]).

For further discussion we introduce the following notation ---
consider the indices as a (finite) set $\D$. Given a function
$\dsize{{\g_1\ldots \g_k}\choose {\d_1 \ldots  \d_k}}$  
%$\pmatrix \g_1\dots \g_k\\\d_1\dots \d_k\endpmatrix$ 
from the marked
points of a $k$-box to $\D$, we define the corresponding
{\it basic tensor} to be
\[ 
T^{j_1\dots j_k}_{i_1\dots i_k}= \left\{
		\begin{array}{ll}
		1 & {\mbox{if}}\ \ i_1=\d_1, \ i_2=\d_2 {\mbox{ \ etc.}}\\ 
		{} & {\mbox{and}} \ \ j_1=\g_1, \ j_2=\g_2 {\mbox{ \ etc.}}\\
		0 & {\mbox{otherwise.}}
		\end{array}
	\right.
\]

%%%%%%%%%%%%%%%%%%%%%%%%%%%%%%%%%%%%%%%%
%\[
%T^{j_1\dots j_k}_{i_1\dots i_k}= \cases
%1 & {\text{if}} \ \ i_1=\d_1, \ i_2=\d_2 {\text{ \ etc.}}\\
%{} & {\text{and}} \ \ j_1=\g_1, \ j_2=\g_2 {\text{ \ etc.}}\\
% 0 & {\text{otherwise.}}\endcases
%\]
%%%%%%%%%%%%%%%%%%%%%%%%%%%%%%%%%%%%%%
If $\frak S(\D)$ is the free semigroup on $\D$, $\p T$ is then the
word \newline $\g_1\g_2\g_3\dots \g_k\d_k\d_{k-1}\dots\d_1$ in $\frak S(\D)$, and we will use
the notation
\fbox{\parbox{2.3cm}{
$\dsize{
\begin{array}{c}
\g_1\ \g_2\ldots \g_k \\
\d_1\ \d_2 \ldots  \d_k
\end{array}
	}$
}} 
%%%%%%%%%%%%%%%%%%%%%%%%%%%
%$\boxed{\matrix \g_1\g_2\dots \g_k\\ \d_1\d_2\dots
%\d_k\endmatrix}$
%%%%%%%%%%%%%%%%%%%%%%%%%
for this basic tensor.

The planar algebra $P^\otimes$ is not terribly interesting by
itself (and there seems  to be no reason to limit the contractions
allowed to planar ones). But one may look for planar subalgebras.
One way is to take a set $\{A_i\in P^\otimes_{k_i}\}$ and look at
the planar subalgebra $P_k(A_i)$ they generate. The calculation of
$P_k(A_i)$ as a function of the $A_i$'s can be extremely difficult.
While it is easy enough to decide if the $A_i$'s are in the TL
subalgebra, we will see in the next example that the question of
whether $P_k(A_i)\neq P_k^{\otimes}$ is undecidable, even for $k=1$.

Note that if the tensors $A_i$ have only 0--1 entries, the
partition function will be simply the number of ``edge colourings"
of the network by $n$ colours with colourings allowed only if they
correspond to a non-zero entry of the tensor label at each box.

The next example gives a situation where we {\it can} say
$P_k(A_i)\neq P^{\otimes}_k$.

\bigskip{\bf Example 2.7: Finitely generated groups.}
As in 2.6, if $\D$ is a set, $\frak S(\D)$ will denote the free
semigroup on $\D$, and $F(\D)$ will denote the free group on $\D$.
We define the map
alt:$\frak S(\D)\to F(\D)$ by alt$(\g_1\dots\g_m)=
\g_1\g_2^{-1}\g_3\g_4^{-1}\dots \g_m^{\pm 1}$,
(where the $+$ sign occurs if $m$ is odd, -- if $m$ is
even). Note that alt is only a homomorphism from the subsemigroup
of words of even length.

Now let $\G$ be a discrete group and $\D$ a finite set, together
with a function $\d\mapsto\tilde\d$ from $\D$ to $\G$. There is
then a natural map $\phi: F(\D)\to\G$ defined by
$\phi(\d)=\tilde\d$. Let
$V$ be the vector space with basis $\D$. Use $V$ and the basis $\D$
to form the planar algebra $P^{\otimes}$ of $\S$2.6. Recall that,
for a basic tensor
$T\in P^{\otimes}_k$, \ $\p T$ is the element of $\frak S(\D)$
obtained by reading around the boundary of $T$. Let $*$ denote the
involution on $\frak S(\D)$ given by writing words backwards.

\bigskip{\bf Definition.} $P^{\G,\D}=\cup_k P_k^{\G,\D}$ is the
linear span of all basic tensors $T$ such that
$\phi({\op{alt}}(\p T))=1$ in $\G$.

\ni\bs{\bf Proposition 2.7.1} $P^{\G,\D}$ is a planar $*$-subalgebra
of $P^{\otimes}$.

\bs
{\bf Proof}. By 1.18, it suffices to show $P^{\G,\D}$ is a unital
subalgebra invariant under the annular category ${\cal A} (\emptyset)$. If $T$ is a basic
tensor in $P_k^{\G,\D}$, let $\p^+T$ (resp.~$\p^-T$) be the element
of $\frak S(\D)$ obtained by reading along the top of $T$
(resp.~the bottom), so $\p T=\p^+ T(\p^-T)^*$. Then for the
product $T_1T_2$ to be non-zero, $\p^-T_1=\p^+T_2$.
In the product alt$(T_1)$alt$(T_2)$, the last letter of
$(\p^-T_1)^*$ is then the same as the first letter of $\p T_2$, but
with opposite sign. Thus the contribution of $\p^-T_1$ cancels with
that of $\p^+T_2$ and $\phi({\op{alt}}(T_1){\op{alt}}(T_2))=1$.
So $P^{\G,\D}$ is a subalgebra, clearly unital and self-adjoint.

Now consider a typical ${\cal A} (\emptyset)$ element $C$ applied to a basic tensor
$T$ as in Figure~2.7.2. 
%%%%%%%%%%%%%%%%%%%%%%%%%%%%%%%%%%%%
\[
	\begin{picture}(0,0)%
\epsfig{file=ps/pic30.ps}%
\end{picture}%
\setlength{\unitlength}{0.00041700in}%
%
\begingroup\makeatletter\ifx\SetFigFont\undefined
% extract first six characters in \fmtname
\def\x#1#2#3#4#5#6#7\relax{\def\x{#1#2#3#4#5#6}}%
\expandafter\x\fmtname xxxxxx\relax \def\y{splain}%
\ifx\x\y   % LaTeX or SliTeX?
\gdef\SetFigFont#1#2#3{%
  \ifnum #1<17\tiny\else \ifnum #1<20\small\else
  \ifnum #1<24\normalsize\else \ifnum #1<29\large\else
  \ifnum #1<34\Large\else \ifnum #1<41\LARGE\else
     \huge\fi\fi\fi\fi\fi\fi
  \csname #3\endcsname}%
\else
\gdef\SetFigFont#1#2#3{\begingroup
  \count@#1\relax \ifnum 25<\count@\count@25\fi
  \def\x{\endgroup\@setsize\SetFigFont{#2pt}}%
  \expandafter\x
    \csname \romannumeral\the\count@ pt\expandafter\endcsname
    \csname @\romannumeral\the\count@ pt\endcsname
  \csname #3\endcsname}%
\fi
\fi\endgroup
\begin{picture}(4844,3944)(879,-3683)
\put(2326,-1411){\makebox(0,0)[lb]{\smash{\SetFigFont{12}{14.4}{rm}$\gamma_1$}}}
\put(3226,-1411){\makebox(0,0)[lb]{\smash{\SetFigFont{12}{14.4}{rm}$\gamma_2$}}}
\put(4051,-1411){\makebox(0,0)[lb]{\smash{\SetFigFont{12}{14.4}{rm}$\gamma_3$}}}
\put(2251,-2161){\makebox(0,0)[lb]{\smash{\SetFigFont{12}{14.4}{rm}$\gamma_6$}}}
\put(3076,-2236){\makebox(0,0)[lb]{\smash{\SetFigFont{12}{14.4}{rm}$\gamma_5$}}}
\put(4051,-2236){\makebox(0,0)[lb]{\smash{\SetFigFont{12}{14.4}{rm}$\gamma_4$}}}
\end{picture}

\]
\begin{center}
 	Figure 2.7.2
\end{center}
%%%%%%%%%%%%%%%%%%%%%%%%%%%%%%%%%%%

\noindent This tensor is a sum of basic tensors, the sum ranging
over all functions from the curves in the diagram to $\D$. If $R$
is a particular basic tensor in the sum, notice that the non through-strings in $C$
contribute to alt$(\p R)$ in two ways --- either they occur in
cancelling pairs or, if their beginning and end are separated by
the left-hand side of the diagram, they change alt$(\p R)$ by
conjugation (eliminate all the cancelling pairs first to see this).
Thus the conjugacy class in $F(\D)$ of alt$(R)$ (and alt$(T)$) is
not changed by removing all non through-strings. Once this is done,
however, the word around the outer boundary is just an even cyclic
permutation of the word $w'$ around the inner boundary, so
$\phi({\op{alt}} \ w)=1\Longleftrightarrow \phi({\op{alt}}
(w'))=1$. Thus $\P^{\G,\D}$ is invariant under ${\cal A} (\emptyset)$. \qed

\bigskip
Note that the basic tensors in $P^{\G,\D}$ are unchanged if we
change $\sim:\D\to\G$ by right multiplication by an element of
$\G$. So we may suppose there is an element $e$ of $\D$ with
$\tilde e=1\in\G$. To denote this situation we will say simply
``$e\in\D$".

Let $G\!\subseteq\G$ (resp.~$G'$) $=\{\phi({\op{alt}}(\p^+ T))\mid T
{\mbox{ \ a\  basic\  tensor in }} P^{\G,\D}_{2k} \
({\text{resp.}} \ P^{\G,\D}_k), \ k\!\in\!\Bbb N\}$.

\bs\ni
{\bf Lemma 2.7.3} $G$ is the subgroup
$\<\tilde\D\tilde\D^{-1}\>$ of $\G$ generated by
$\tilde\D\tilde\D^{-1}$, and if
$e\in\D$, \ $G=G'$.

\bs
{\sc Proof}. The definition of alt implies immediately that
$G\subseteq\<\tilde\D\tilde\D^{-1}\>$. That $G=G^{-1}$ follows
from Figure~2.7.4
%%%%%%%%%%%%%%%%%%%%%%%%%%%
\[
        \begin{picture}(0,0)%
\epsfig{file=ps/pic31.ps}%
\end{picture}%
\setlength{\unitlength}{0.00041700in}%
%
\begingroup\makeatletter\ifx\SetFigFont\undefined
% extract first six characters in \fmtname
\def\x#1#2#3#4#5#6#7\relax{\def\x{#1#2#3#4#5#6}}%
\expandafter\x\fmtname xxxxxx\relax \def\y{splain}%
\ifx\x\y   % LaTeX or SliTeX?
\gdef\SetFigFont#1#2#3{%
  \ifnum #1<17\tiny\else \ifnum #1<20\small\else
  \ifnum #1<24\normalsize\else \ifnum #1<29\large\else
  \ifnum #1<34\Large\else \ifnum #1<41\LARGE\else
     \huge\fi\fi\fi\fi\fi\fi
  \csname #3\endcsname}%
\else
\gdef\SetFigFont#1#2#3{\begingroup
  \count@#1\relax \ifnum 25<\count@\count@25\fi
  \def\x{\endgroup\@setsize\SetFigFont{#2pt}}%
  \expandafter\x
    \csname \romannumeral\the\count@ pt\expandafter\endcsname
    \csname @\romannumeral\the\count@ pt\endcsname
  \csname #3\endcsname}%
\fi
\fi\endgroup
\begin{picture}(7349,5310)(1165,-5383)
\put(1877,-2762){\makebox(0,0)[lb]{\smash{\SetFigFont{12}{14.4}{rm}$T$}}}
\put(6376,-2761){\makebox(0,0)[lb]{\smash{\SetFigFont{12}{14.4}{rm}$T$}}}
\put(1201,-1561){\makebox(0,0)[lb]{\smash{\SetFigFont{12}{14.4}{rm}$\phi ($alt$(\partial^+T))=w$}}}
\put(1201,-4186){\makebox(0,0)[lb]{\smash{\SetFigFont{12}{14.4}{rm}$\phi ($alt$(\partial^+T))=w^{-1}$}}}
\put(5701,-361){\makebox(0,0)[lb]{\smash{\SetFigFont{12}{14.4}{rm}$\phi ($alt$(\partial^+T))=w^{-1}$}}}
\put(5476,-5311){\makebox(0,0)[lb]{\smash{\SetFigFont{12}{14.4}{rm}$\phi ($alt$(\partial^+T))=w$}}}
\end{picture}

\]
\begin{center}
	Figure 2.7.4
\end{center}
%%%%%%%%%%%%%%%%%%%%%%%%%%%%%
\noindent That $G$ is a group follows
from Figure~2.7.5 
%%%%%%%%%%%%%%%%%%%%%%%%%%%
\[
	\begin{picture}(0,0)%
\epsfig{file=ps/pic32.ps}%
\end{picture}%
\setlength{\unitlength}{0.00041700in}%
%
\begingroup\makeatletter\ifx\SetFigFont\undefined
% extract first six characters in \fmtname
\def\x#1#2#3#4#5#6#7\relax{\def\x{#1#2#3#4#5#6}}%
\expandafter\x\fmtname xxxxxx\relax \def\y{splain}%
\ifx\x\y   % LaTeX or SliTeX?
\gdef\SetFigFont#1#2#3{%
  \ifnum #1<17\tiny\else \ifnum #1<20\small\else
  \ifnum #1<24\normalsize\else \ifnum #1<29\large\else
  \ifnum #1<34\Large\else \ifnum #1<41\LARGE\else
     \huge\fi\fi\fi\fi\fi\fi
  \csname #3\endcsname}%
\else
\gdef\SetFigFont#1#2#3{\begingroup
  \count@#1\relax \ifnum 25<\count@\count@25\fi
  \def\x{\endgroup\@setsize\SetFigFont{#2pt}}%
  \expandafter\x
    \csname \romannumeral\the\count@ pt\expandafter\endcsname
    \csname @\romannumeral\the\count@ pt\endcsname
  \csname #3\endcsname}%
\fi
\fi\endgroup
\begin{picture}(5744,3236)(1179,-3384)
\put(2551,-2536){\makebox(0,0)[lb]{\smash{\SetFigFont{12}{14.4}{rm}$T$}}}
\put(5701,-2536){\makebox(0,0)[lb]{\smash{\SetFigFont{12}{14.4}{rm}$R$}}}
\put(1576,-1261){\makebox(0,0)[lb]{\smash{\SetFigFont{12}{14.4}{rm}$\phi($alt$(\partial^+T))$}}}
\put(5026,-1261){\makebox(0,0)[lb]{\smash{\SetFigFont{12}{14.4}{rm}$\phi$(alt$(\partial^+R))$}}}
\put(2626,-436){\makebox(0,0)[lb]{\smash{\SetFigFont{12}{14.4}{rm}$\phi($alt$(\partial^+T))\phi($alt$(\partial^+R))$}}}
\end{picture}

\]
\begin{center}
	Figure 2.7.5
\end{center}
%%%%%%%%%%%%%%%%%%%%%%%%%%%%%
\noindent Also $G$ contains $\tilde\D\tilde\D^{-1}$ since
$\bsp$ is in $P_2^{\G,\D}$ where
for $\g\in\D$, \ $\boxa{\g}$ is the ``diagonal" tensor
$\blong$. That $G=G'$ if $e\in\D$ is
easily seen by attaching $\boxb{e}$ to the right of basic tensors
in $P_k^{\G,\D}$ when $k$ is even. \qed

\bigskip
We see that, if $e\in\D$, a basis for $P^{\G,\D}$ is formed by all
random walks on $G$, starting and ending at $1\in\G$, where the odd
transitions correspond to multiplying by a $\tilde \d$ for each $\d\in\D$,
and the even ones by $\tilde\d^{-1}$ for $\d\in\D$. If
$\tilde\D=\tilde\D^{-1}$ and $\tilde{}$ is injective, these are
just random walks on the Cayley graph of $G$.

If $e\in\D$, each basic tensor $T\in P^{\G,\D}$ gives the {\it
relation} alt$(\p T)$ in $G$, thinking of $G$ as being presented on
$\D\backslash\{e\}$.


